\documentclass[11pt]{article}
\usepackage[T1]{fontenc}
\usepackage[utf8]{inputenc}
\usepackage{lmodern}
\usepackage[a4paper,margin=27mm]{geometry}
\usepackage{amsmath,amssymb,amsthm,mathtools}
\usepackage{bm}
\usepackage{booktabs,longtable,array}
\usepackage{graphicx}
\usepackage{microtype}
\usepackage{xcolor}
\usepackage{enumitem}
\usepackage{url}
\usepackage{hyperref}
\usepackage[nameinlink,noabbrev]{cleveref}
\usepackage{fancyhdr}
\usepackage{listings}

% ed.: corpus_meta.tex was not shipped in the archive; the input now
% reads a reconstructed stub defining \CorpusFileCount, \CorpusCommit,
% and \CorpusDigest as placeholders (see the note in that file).
% ed.: this file was referenced by % ed.: this file was referenced by % ed.: this file was referenced by \input{corpus_meta.tex} but not
% shipped in the archive (the README lists it, together with
% corpus_inventory.tex and corpus_audit.txt, as part of the intended
% package).  The three macros below are reconstructed as placeholders
% so the document compiles; the actual snapshot metadata of the
% generator's corpus audit is unrecoverable.
\newcommand{\CorpusFileCount}{[not shipped]}
\newcommand{\CorpusCommit}{\texttt{[not shipped]}}
\newcommand{\CorpusDigest}{\texttt{[not shipped]}}
 but not
% shipped in the archive (the README lists it, together with
% corpus_inventory.tex and corpus_audit.txt, as part of the intended
% package).  The three macros below are reconstructed as placeholders
% so the document compiles; the actual snapshot metadata of the
% generator's corpus audit is unrecoverable.
\newcommand{\CorpusFileCount}{[not shipped]}
\newcommand{\CorpusCommit}{\texttt{[not shipped]}}
\newcommand{\CorpusDigest}{\texttt{[not shipped]}}
 but not
% shipped in the archive (the README lists it, together with
% corpus_inventory.tex and corpus_audit.txt, as part of the intended
% package).  The three macros below are reconstructed as placeholders
% so the document compiles; the actual snapshot metadata of the
% generator's corpus audit is unrecoverable.
\newcommand{\CorpusFileCount}{[not shipped]}
\newcommand{\CorpusCommit}{\texttt{[not shipped]}}
\newcommand{\CorpusDigest}{\texttt{[not shipped]}}


\hypersetup{
  colorlinks=true,
  linkcolor=blue!45!black,
  citecolor=green!35!black,
  urlcolor=blue!55!black,
  pdftitle={Dyadic-Comb Moments of the Fabius and Rvachev Functions},
  pdfauthor={OpenAI GPT-5.6 Pro, prepared for Vladimir Reshetnikov}
}

\pagestyle{fancy}
\fancyhf{}
\lhead{Dyadic-comb Fabius/Rvachev sums}
\rhead{\thepage}
\setlength{\headheight}{14pt}

\newtheorem{theorem}{Theorem}[section]
\newtheorem{lemma}[theorem]{Lemma}
\newtheorem{corollary}[theorem]{Corollary}
\newtheorem{proposition}[theorem]{Proposition}
\newtheorem{conjecture}[theorem]{Conjecture}
\theoremstyle{definition}
\newtheorem{definition}[theorem]{Definition}
\newtheorem{example}[theorem]{Example}
\theoremstyle{remark}
\newtheorem{remark}[theorem]{Remark}

\newcommand{\up}{\operatorname{up}}
\newcommand{\sinc}{\operatorname{sinc}}
\newcommand{\Li}{\operatorname{Li}}
\newcommand{\e}{\mathrm e}
\newcommand{\ii}{\mathrm i}
\newcommand{\N}{\mathbb N}
\newcommand{\Z}{\mathbb Z}
\newcommand{\C}{\mathbb C}
\newcommand{\R}{\mathbb R}
\newcommand{\fall}[2]{#1^{\underline{#2}}}
\newcommand{\tm}{\tau}
\newcommand{\dd}{\,\mathrm d}
\newcommand{\D}{\mathcal D}
\newcommand{\J}{\mathcal J}
\newcommand{\B}{\mathcal B}
\newcommand{\Rema}{\mathcal R}

\title{\textbf{Dyadic-Comb Moments of the Fabius and Rvachev Functions}\\[4pt]
\large Exact rational formulae, sinc-zero superconvergence, spectral Dirichlet values, and iterated sums}
\author{Research report prepared for Vladimir Reshetnikov}
\date{28 August 2026}

\begin{document}
\maketitle

\begin{abstract}
Let $F$ be the Fabius function on $[0,1]$, let $\up$ be Rvachev's compactly
supported bump, and put $N=2^m$.  This report studies the weighted dyadic-comb
rule
\[
 S_{m,p}=\frac1N\sum_{k=0}^{N}\left(\frac{k}{N}\right)^pF\!\left(\frac{k}{N}\right)
\]
and its extensions to complex powers and repeated summation.  The first main
result is a finite Thue--Morse/Appell formula for every dyadic value and hence a
closed rational expression for $S_{m,p}$ at arbitrary $m,p\in\N_0$.  Its
generating function is rational after negative-index polylogarithms are
expanded, and an $O(m2^m)$ sparse-difference algorithm replaces the apparent
quadratic convolution.

The second main result explains an unexpectedly strong phenomenon.  The
Fourier transform
\[
 H(z)=\widehat{\up}(z)=\prod_{\nu\ge1}\sinc\!\left(\frac{z}{2^\nu}\right)
\]
has a zero of multiplicity $m+1+v_2(\ell)$ at every nonzero alias
$2\pi 2^m\ell$.  Consequently the full Rvachev comb integrates every
polynomial of degree at most $m$ exactly.  For the Fabius half-bump this gives,
for $m\ge p+1$,
\[
 S_{m,p}=I_p+\frac{1}{2N}
 +\sum_{q=1}^{\lfloor(p+1)/2\rfloor}
   \frac{B_{2q}}{(2q)!}\fall{p}{2q-1}N^{-2q},
 \qquad I_p=\int_0^1x^pF(x)\dd x,
\]
with no unshown remainder.  The boundary case $m=p$ is also exact when $p$ is
even.  A first-alias calculation produces a new family of exact spectral
Dirichlet evaluations; for example
\[
 \sum_{q\ge1\atop q\ \mathrm{odd}}\frac{H(\pi q)}{q^2}=\frac{\pi^2}{18},
 \qquad
 \sum_{q\ge1\atop q\ \mathrm{odd}}\frac{H(\pi q)}{q^4}=\frac{23\pi^4}{4050}.
\]

For arbitrary $\alpha\in\C$, both the continuous Mellin moment and the finite
comb are entire functions of $\alpha$.  Integer truncation is replaced by a
universal Euler--Maclaurin expansion with falling factorials.  Finally, an
inclusive $n$-fold sum is reduced exactly to $n$ ordinary comb moments by a
binomial kernel; its continuum limit is a beta-weighted Fabius moment.  The
report includes exact checks, high-precision experiments, conjectures about
alias sharpness and spectral positivity, and a recursive inventory of the
repository corpus used to avoid duplicating its existing integral,
asymptotic, interpolation, and $q$-series developments.
\end{abstract}

\tableofcontents

\section{Scope, corpus audit, and novelty boundary}

The source corpus was taken recursively from
\begin{center}
\url{https://github.com/VladimirReshetnikov/ProveIt/tree/main/Analysis/FabiusFunction/docs}.
\end{center}
The checked snapshot contains \CorpusFileCount{} \texttt{.tex} files at Git
commit \CorpusCommit; the combined content digest begins \CorpusDigest.  The
machine-readable audit is included as \texttt{corpus\_audit.txt}, and the full
relative-path inventory appears in \cref{app:corpus}.

The corpus already develops a substantial common foundation: exact dyadic
values, Thue--Morse structure, $q$-Pochhammer and $q$-binomial forms, sinc
products and Fourier transforms, continuous moments, Bell/Bernoulli
polynomials, endpoint asymptotics involving Lambert $W$, inverse-Fabius
questions, interpolation, Richardson extrapolation, and several classes of
integrals.  Those ingredients are used rather than re-presented as purported
new discoveries.  The emphasis here is deliberately narrower and discrete:

\begin{enumerate}[label=(\roman*),leftmargin=2.4em]
\item a closed generating object for \emph{weighted comb sums}, not merely a
      formula for an individual dyadic value;
\item exactness and defect formulae obtained from the \emph{multiplicity} of
      dyadic sinc zeros under Poisson summation;
\item spectral Dirichlet values extracted from the first nonzero alias;
\item complex powers, including negative powers, treated through an entire
      Mellin parameter; and
\item inclusive, strict, and fractional-order iterated sums.
\end{enumerate}

A repository audit can rule out literal and close structural duplication in a
finite corpus; it cannot certify priority against all unpublished mathematics.
Accordingly, statements proved below are labeled theorems, computational
observations are labeled as such, and wider arithmetic patterns are explicitly
presented as conjectures.

\section{Normalization and basic identities}

\subsection{Rvachev's bump and the Fabius half-bump}

Throughout,
\[
 \sinc z=\begin{cases}\sin z/z,&z\ne0,\\1,&z=0,\end{cases}
 \qquad
 H(z)=\prod_{\nu=1}^{\infty}\sinc\!\left(\frac{z}{2^\nu}\right).
\]
With the Fourier convention
\(
 \widehat f(z)=\int_{\R}f(x)\e^{-\ii zx}\dd x,
\)
Rvachev's function $u=\up$ is the even $C^\infty$ function supported on
$[-1,1]$ satisfying $\widehat u=H$, $u(0)=1$, and $\int u=1$.  On $[0,1]$ we
use
\begin{equation}\label{eq:F-up-relation}
 F(x)=u(x-1)=u(1-x).
\end{equation}
Thus
\begin{equation}\label{eq:partition}
 F(x)+F(1-x)=1,
 \qquad
 u(x)+u(1-x)=1
 \quad(0\le x\le1).
\end{equation}
Both $F$ at $0$ and $1-F$ at $1$ are flat: every one-sided derivative
vanishes.  Equivalently,
\begin{equation}\label{eq:flat}
 F(x)=O(x^A),\qquad 1-F(1-x)=O(x^A)
 \quad(x\downarrow0)
\end{equation}
for every $A>0$.  This flatness is what makes negative powers harmless later.

\subsection{A probabilistic model}

Let $U_1,U_2,\ldots$ be independent uniform random variables on $[-1,1]$, and
put
\begin{equation}\label{eq:X}
 X=\sum_{\nu=1}^{\infty}2^{-\nu}U_\nu.
\end{equation}
Then $X$ has density $u$, because its characteristic function is $H$.  Its
moment-generating function is
\begin{equation}\label{eq:M}
 M(t)=\mathbb E\e^{tX}
 =\prod_{\nu=1}^{\infty}
   \frac{\sinh(t/2^\nu)}{t/2^\nu}.
\end{equation}
Write $\mu_r=\mathbb E X^r$.  Symmetry gives $\mu_{2r+1}=0$.  From
\[
 \log\frac{\sinh z}{z}
 =\sum_{r\ge1}\frac{2^{2r-1}B_{2r}}{r(2r)!}z^{2r}
\]
we obtain the even cumulants
\begin{equation}\label{eq:cumulants}
 \kappa_{2r}=\frac{2^{2r-1}B_{2r}}{r(2^{2r}-1)},
 \qquad \kappa_{2r+1}=0,
\end{equation}
and hence
\begin{equation}\label{eq:mu-bell}
 \mu_n=Y_n(0,\kappa_2,0,\kappa_4,\ldots,\kappa_n),
\end{equation}
where $Y_n$ is the complete exponential Bell polynomial.  In particular all
$\mu_n$ are rational.  The first values are $\mu_0=1$, $\mu_2=1/9$, and
$\mu_4=19/675$.

Define the Appell moment polynomials
\begin{equation}\label{eq:P-m}
 P_m(y)=\mathbb E(y-X)^m
 =\sum_{r=0}^{m}(-1)^r\binom mr\mu_r y^{m-r}
 =m![t^m]\e^{yt}M(t).
\end{equation}
Only the even $r$ survive, so the signs disappear in the explicit sum.

\section{Continuous moments in closed form}

The continuous moments are needed both as limiting integrals and as constants
in the exact comb formula.

\begin{theorem}[Exponential generating function]\label{thm:moment-egf}
Let
\(
 I_p=\int_0^1x^pF(x)\dd x.
\)
Then
\begin{equation}\label{eq:I-egf}
 \sum_{p=0}^{\infty}I_p\frac{t^p}{p!}
 =\frac{\e^t}{t}-\frac{\e^tM(t)}{\e^t-1}.
\end{equation}
Consequently
\begin{equation}\label{eq:I-closed}
 I_p=\frac1{p+1}\left[
 1-\sum_{j=0}^{p+1}\binom{p+1}{j}\mu_j B_{p+1-j}(1)
 \right],
\end{equation}
where $B_n(x)$ is the Bernoulli polynomial.  In particular, $I_p\in\mathbb Q$
for every $p\in\N_0$.
\end{theorem}

\begin{proof}
Set $A(t)=\int_0^1u(x)\e^{tx}\dd x$.  Evenness gives
$A(t)+A(-t)=M(t)$, while the partition identity \eqref{eq:partition} gives
\[
 A(t)+\e^tA(-t)=\frac{\e^t-1}{t}.
\]
Solving the two equations yields
\[
 A(-t)=\frac1t-\frac{M(t)}{\e^t-1}.
\]
By \eqref{eq:F-up-relation},
\[
 \int_0^1F(x)\e^{tx}\dd x
 =\e^tA(-t),
\]
which proves \eqref{eq:I-egf}.  Expanding
\(
 t\e^t/(\e^t-1)=\sum_{n\ge0}B_n(1)t^n/n!
\)
and using $M(t)=\sum\mu_jt^j/j!$ gives \eqref{eq:I-closed}.
\end{proof}

The same calculation gives the positive-half moments of $u$:
\begin{equation}\label{eq:a-r}
 a_r:=\int_0^1x^ru(x)\dd x
 =\frac1{r+1}\sum_{j=0}^{r+1}
   \binom{r+1}{j}\mu_jB_{r+1-j}(1)
 =\frac1{r+1}-I_r.
\end{equation}
These formulas connect the Bell-polynomial moment algebra to the Bernoulli
polynomials without invoking numerical quadrature.

\section{A finite Thue--Morse formula for every dyadic value}

Let
\[
 N=2^m,\qquad h=N^{-1},\qquad
 \tm(a)=(-1)^{s_2(a)},
\]
where $s_2(a)$ is the binary digit sum.

\begin{theorem}[Dyadic Appell--Thue--Morse formula]\label{thm:dyadic-values}
For $m\in\N_0$ and $0\le k\le N$,
\begin{equation}\label{eq:dyadic-value}
 F\!\left(\frac{k}{N}\right)
 =\frac{2^{-m(m+1)/2}}{m!}
   \sum_{a=0}^{k-1}\tm(a)P_m\bigl(2(k-a)-1\bigr),
\end{equation}
where the empty sum for $k=0$ is zero.
\end{theorem}

\begin{proof}
Put $n=m+1$ and split the random series \eqref{eq:X} after $n$ terms.  After
multiplication by $2^n$,
\[
 2^nX=\sum_{\nu=1}^{n}2^{n-\nu}U_\nu+X',
\]
where $X'$ is an independent copy of $X$.  The density of the finite sum on the
right is a nonuniform Irwin--Hall spline.  Since its half-widths are
$1,2,4,\ldots,2^{n-1}$, their subset sums are exactly the integers
$0,\ldots,2^n-1$, and inclusion--exclusion gives
\[
 f_n(y)=\frac{2^{-n(n+1)/2}}{(n-1)!}
 \sum_{a=0}^{2^n-1}\tm(a)
 (y+2^n-1-2a)_+^{n-1}.
\]
The density scaling identity is
\[
 2^{-n}u(y/2^n)=\mathbb E f_n(y-X').
\]
At $x=k/N-1$ one has $y=2(k-N)$, and the truncated-power argument becomes
$2(k-a)-1-X'$.  Since $X'\in[-1,1]$, it is nonnegative exactly when
$a\le k-1$ (endpoint equalities are immaterial).  Taking the expectation
produces $P_m(2(k-a)-1)$ and yields \eqref{eq:dyadic-value}.
\end{proof}

\begin{example}
For $m=2$, $P_2(y)=y^2+1/9$ and the prefactor is $1/16$.  Thus
\[
 F(1/4)=\frac1{16}\left(1+\frac19\right)=\frac5{72},
 \qquad F(1/2)=\frac12.
\]
For $m=3$, $P_3(y)=y^3+y/3$, and one obtains
\[
 \bigl(F(k/8)\bigr)_{k=0}^8
 =\left(0,\frac1{288},\frac5{72},\frac{73}{288},\frac12,
 \frac{215}{288},\frac{67}{72},\frac{287}{288},1\right).
\]
\end{example}

\section{Closed forms for polynomially weighted dyadic combs}

\subsection{A direct finite formula}

For $p\in\N_0$ define the right-endpoint comb
\begin{equation}\label{eq:Smp}
 S_{m,p}=h\sum_{k=0}^{N}x_k^pF(x_k),
 \qquad x_k=kh.
\end{equation}
The $k=0$ term is zero, including at $p=0$ under the natural convention.
Substitution of \eqref{eq:dyadic-value} immediately gives the requested finite
closed form.

\begin{corollary}[Finite rational comb formula]\label{cor:finite-comb}
For every $m,p\in\N_0$,
\begin{equation}\label{eq:finite-comb}
 S_{m,p}=
 \frac{2^{-m(m+1)/2}}{m!N^{p+1}}
 \sum_{k=1}^{N}k^p
 \sum_{a=0}^{k-1}\tm(a)P_m\bigl(2(k-a)-1\bigr).
\end{equation}
In particular $S_{m,p}\in\mathbb Q$.
\end{corollary}

After putting $d=k-a$, the inner power sums can be written with Bernoulli
polynomials:
\[
 \sum_{d=1}^{L}d^r
 =\frac{B_{r+1}(L+1)-B_{r+1}(1)}{r+1}.
\]
Expanding $(a+d)^pP_m(2d-1)$ therefore turns \eqref{eq:finite-comb} into a
single finite Thue--Morse sum of explicit Bernoulli polynomials.  This is useful
for symbolic elimination; the generating function below is more compact.

\subsection{Generating function and an $O(m2^m)$ algorithm}

Define
\begin{align}
 \Theta_m(z)&=\sum_{a=0}^{N-1}\tm(a)z^a
             =\prod_{j=0}^{m-1}(1-z^{2^j}),\label{eq:theta}\\
 B_m^\star(z)&=\sum_{d\ge1}P_m(2d-1)z^d,\label{eq:Bstar}\\
 G_m(z)&=\frac{2^{-m(m+1)/2}}{m!}\Theta_m(z)B_m^\star(z).\label{eq:Gm}
\end{align}

\begin{theorem}[Dyadic-value generating function]\label{thm:Gm}
The coefficient of $z^k$ in $G_m(z)$ equals $F(k/N)$ for $0\le k\le N$ and
equals $1$ for every $k\ge N$.  Hence
\begin{equation}\label{eq:Am}
 A_m(z):=G_m(z)-\frac{z^{N+1}}{1-z}
       =\sum_{k=0}^{N}F(k/N)z^k
\end{equation}
is a polynomial, and
\begin{equation}\label{eq:log-derivative-comb}
 \boxed{
 S_{m,p}=N^{-p-1}
 \left.\left(z\frac{\dd}{\dd z}\right)^pA_m(z)\right|_{z=1}.}
\end{equation}
All apparent singularities at $z=1$ in this representation are removable.
\end{theorem}

\begin{proof}
The coefficient convolution in \eqref{eq:Gm} is precisely
\eqref{eq:dyadic-value}.  Once $k\ge N$, all $a=0,\ldots,N-1$ occur.  The
Prouhet cancellation
\[
 \sum_{a=0}^{N-1}\tm(a)a^r=0\quad(r<m),
 \qquad
 \sum_{a=0}^{N-1}\tm(a)a^m
 =(-1)^m m!2^{m(m-1)/2}
\]
shows that only the leading term of the monic polynomial $P_m$ survives, and
its value after the factor $2^m$ from $P_m(2(k-a)-1)$ is exactly the reciprocal
of the prefactor.  Thus the tail coefficients are $1$.  Applying the Euler
operator $z\dd/\dd z$ multiplies the $k$th coefficient by $k$, proving
\eqref{eq:log-derivative-comb}.
\end{proof}

The series $B_m^\star$ is rational.  Indeed,
\begin{equation}\label{eq:B-polylog}
 B_m^\star(z)=
 \sum_{r=0}^{\lfloor m/2\rfloor}\binom m{2r}\mu_{2r}
 \sum_{j=0}^{m-2r}\binom{m-2r}{j}2^j(-1)^{m-2r-j}\Li_{-j}(z),
\end{equation}
where every negative-index polylogarithm is rational:
\(
 \Li_{-j}(z)=zA_j(z)/(1-z)^{j+1}
\)
with $A_j$ an Eulerian polynomial.  Equations
\eqref{eq:Am}--\eqref{eq:B-polylog} are therefore a finite algebraic closed
form for every natural $p$.

Computationally, one should not form the double sum.  Initialize the coefficient
array with $P_m(2d-1)$ and multiply successively by
$(1-z),(1-z^2),\ldots,(1-z^{2^{m-1}})$.  Each multiplication is one sparse
backward difference, so all $N+1$ dyadic values cost $O(mN)$ arithmetic
operations and $O(N)$ storage.  The accompanying \texttt{experiments.py}
implements this exact algorithm with rational arithmetic.

\section{Poisson summation and dyadic exactness}

\subsection{Exact polynomial quadrature for the full Rvachev bump}

For a polynomial $P$ define
\[
 Q_m[P]=h\sum_{k=-N}^{N}u(kh)P(kh).
\]
Since $u$ and all of its derivatives vanish at $\pm1$, no endpoint convention
is needed.

\begin{theorem}[Full-bump dyadic exactness]\label{thm:up-exactness}
If $\deg P\le m$, then
\begin{equation}\label{eq:up-exactness}
 Q_m[P]=\int_{-1}^{1}u(x)P(x)\dd x.
\end{equation}
Equivalently, for $0\le p\le m$,
\begin{equation}\label{eq:up-monomial}
 h\sum_{k=-N}^{N}(kh)^pu(kh)=\mu_p.
\end{equation}
\end{theorem}

\begin{proof}
For a monomial, Poisson summation gives
\[
 h\sum_{k\in\Z}(kh)^pu(kh)
 =\sum_{\ell\in\Z}\ii^pH^{(p)}(2\pi N\ell).
\]
At $z_\ell=2\pi2^m\ell\ne0$, each factor
$\sinc(z/2^\nu)$ with
$1\le\nu\le m+1+v_2(\ell)$ vanishes simply.  Thus
\begin{equation}\label{eq:zero-order}
 \operatorname{ord}_{z=z_\ell}H(z)=m+1+v_2(\ell).
\end{equation}
Every nonzero alias therefore vanishes after at most $m$ derivatives.  Only
$\ell=0$ remains, and its value is the moment $\mu_p$.
\end{proof}

This is stronger than rapid convergence: at each level the dyadic comb is an
exact quadrature rule on a growing polynomial space.  For even $p>0$ it also
gives the useful half-grid identity
\begin{equation}\label{eq:up-half-even}
 h\sum_{k=0}^{N}(kh)^pu(kh)=\frac{\mu_p}{2}\qquad(2\mid p,\ 0<p\le m),
\end{equation}
while for $p=0$ the center sample contributes
\(
 h\sum_{k=0}^{N}u(kh)=(1+h)/2.
\)

\subsection{The one-sided Fourier transform}

Let $g_0(x)=F(x)$ on $[0,1]$ and zero outside, with the symmetric half-value at
the jump $x=1$.  A key identity is
\begin{equation}\label{eq:g0hat}
 \widehat g_0(z)
 =-\frac{\e^{-\ii z}}{\ii z}+\frac{H(z)}{\e^{\ii z}-1}.
\end{equation}
To derive it, put
$U_+(z)=\int_0^1u(x)\e^{-\ii zx}\dd x$.  Evenness and
\eqref{eq:partition} give
\[
 U_+(z)+U_+(-z)=H(z),
 \qquad
 U_+(z)+\e^{-\ii z}U_+(-z)=\frac{1-\e^{-\ii z}}{\ii z}.
\]
Solving and shifting by \eqref{eq:F-up-relation} yields
\eqref{eq:g0hat}.  For $g_p(x)=x^pF(x)$,
\begin{equation}\label{eq:gphat}
 \widehat g_p(z)=\ii^p\frac{\dd^p}{\dd z^p}\widehat g_0(z).
\end{equation}

At a nonzero dyadic alias, the quotient
\begin{equation}\label{eq:Rz}
 R(z):=\frac{H(z)}{\e^{\ii z}-1}
\end{equation}
has zero order
\begin{equation}\label{eq:R-order}
 \operatorname{ord}_{z=2\pi N\ell}R(z)=m+v_2(\ell).
\end{equation}
The denominator removes exactly one of the sinc zeros.  This is the source of
the finite rule below.

\subsection{An exact finite Euler--Maclaurin formula}

The right rule \eqref{eq:Smp} differs from the symmetric trapezoidal convention
by $h/2$, because $F(1)=1$ and $F(0)=0$.

\begin{theorem}[Fabius dyadic superconvergence]\label{thm:superconvergence}
Let $m,p\in\N_0$, $N=2^m$, and $I_p$ be given by
\eqref{eq:I-closed}.  If $m\ge p+1$, then
\begin{equation}\label{eq:stabilized}
 \boxed{
 S_{m,p}=I_p+\frac1{2N}
 +\sum_{q=1}^{\lfloor(p+1)/2\rfloor}
   \frac{B_{2q}}{(2q)!}\fall{p}{2q-1}N^{-2q}.}
\end{equation}
The same formula holds in the boundary case $m=p$ whenever $p$ is even.
\end{theorem}

\begin{proof}
Poisson summation, with the half-value at $x=1$, gives
\[
 S_{m,p}-\frac h2
 =I_p+2\Re\sum_{\ell\ge1}\widehat g_p(2\pi N\ell).
\]
By \eqref{eq:R-order}, the contribution from $R(z)$ in
\eqref{eq:g0hat} vanishes whenever $p<m$.  It remains to differentiate the
endpoint term $-\e^{-\ii z}/(\ii z)$.  At $z=2\pi N\ell$,
\[
 \ii^p\frac{\dd^p}{\dd z^p}
 \left(-\frac{\e^{-\ii z}}{\ii z}\right)
 =\sum_{r=0}^{p}\binom pr r!\ii^{1-r}z^{-r-1}.
\]
Pairing $\ell$ and $-\ell$ retains only $r=2q-1$.  Euler's evaluation of
$\zeta(2q)$ converts the result to
\[
 \frac{B_{2q}}{(2q)!}\fall p{2q-1}N^{-2q},
\]
which proves \eqref{eq:stabilized}.  If $p=m$ is even, the leading derivative
of $R$ has odd Fourier parity and its paired real part vanishes.  Alternatively,
combine \cref{thm:up-exactness} with \eqref{eq:up-half-even} and
$u(x)=1-F(x)$.
\end{proof}

For example,
\begin{align*}
 S_{m,0}&=\frac12+\frac1{2N},\\
 S_{m,1}&=\frac{13}{36}+\frac1{2N}+\frac1{12N^2}
 \qquad(m\ge2),\\
 S_{m,2}&=\frac5{18}+\frac1{2N}+\frac1{6N^2}
 \qquad(m\ge2),\\
 S_{m,3}&=\frac{1207}{5400}+\frac1{2N}+\frac1{4N^2}
          -\frac1{120N^4}
 \qquad(m\ge4).
\end{align*}
The formulas are exact identities at the stated dyadic levels, not asymptotic
expansions.

For even $p$ there is a second compact form.  From
\eqref{eq:up-half-even}, for $0<p\le m$ even,
\begin{equation}\label{eq:even-power-form}
 S_{m,p}=\frac1{N^{p+1}}\sum_{k=1}^{N}k^p-\frac{\mu_p}{2}
 =\frac{B_{p+1}(N+1)-B_{p+1}}{(p+1)N^{p+1}}-\frac{\mu_p}{2}.
\end{equation}
This identity makes the rationality and the boundary extension $m=p$ entirely
transparent.

\subsection{All levels: an exact alias remainder}

At levels below the stabilization threshold, the missing term is still
explicit.

\begin{proposition}[Alias remainder]\label{prop:alias-remainder}
For arbitrary $m,p\in\N_0$,
\begin{equation}\label{eq:all-levels}
 S_{m,p}=I_p+\frac1{2N}
 +\sum_{q=1}^{\lfloor(p+1)/2\rfloor}
   \frac{B_{2q}}{(2q)!}\fall p{2q-1}N^{-2q}
 +\Rema_{m,p},
\end{equation}
where the absolutely convergent alias series is
\begin{equation}\label{eq:alias-series}
 \Rema_{m,p}=
 2\Re\sum_{\ell=1}^{\infty}
 \ii^p\left.
 \frac{\dd^p}{\dd z^p}
 \left(\frac{H(z)}{\e^{\ii z}-1}\right)
 \right|_{z=2\pi N\ell}.
\end{equation}
Only valuation classes $v_2(\ell)\le p-m$ can contribute.
\end{proposition}

The last observation follows immediately from \eqref{eq:R-order}.  It converts
an infinite alias set into finitely many $2$-adic valuation strata, each still
containing rapidly decaying odd multiples.  This is a useful starting point
for higher-defect formulae.

\section{The first alias defect and spectral Dirichlet values}

Let
\begin{equation}\label{eq:Ds}
 D(s)=\sum_{q\ge1\atop q\ \mathrm{odd}}\frac{H(\pi q)}{q^s}.
\end{equation}
Because $H(x)$ decreases faster than every inverse power on the real axis,
this Dirichlet series defines an entire function of $s$.

Fix an odd $q$ and put $z_0=2\pi2^m2^rq$, so
$s_0=m+r+1$ sinc factors vanish.  Removing those factors gives the local
expansion
\begin{equation}\label{eq:local-H}
 H(z)=-\frac{H(\pi q)}{z_0^{s_0}}(z-z_0)^{s_0}
       +O\bigl((z-z_0)^{s_0+1}\bigr).
\end{equation}
Since $\e^{\ii z}-1=\ii(z-z_0)+O((z-z_0)^2)$,
\begin{equation}\label{eq:local-R}
 R(z)=\frac{\ii H(\pi q)}{z_0^{s_0}}
       (z-z_0)^{s_0-1}+O\bigl((z-z_0)^{s_0}\bigr).
\end{equation}
At derivative order $p=m$, only the odd aliases ($r=0$) can occur.

\begin{theorem}[First-defect identity]\label{thm:first-defect}
For odd $m\ge1$,
\begin{align}
 &S_{m,m}-I_m-\frac1{2N}
 -\sum_{q=1}^{(m+1)/2}
   \frac{B_{2q}}{(2q)!}\fall m{2q-1}N^{-2q}
 \notag\\
 &\hspace{35mm}=
 \frac{2m!(-1)^{(m+1)/2}}{(2\pi N)^{m+1}}D(m+1).
 \label{eq:first-defect}
\end{align}
Consequently
\begin{equation}\label{eq:D-rational}
 \boxed{\quad D(2r)/\pi^{2r}\in\mathbb Q\qquad(r\ge1).\quad}
\end{equation}
\end{theorem}

\begin{proof}
Insert \eqref{eq:local-R} into \eqref{eq:alias-series} with $p=m$.
The factor $\Re(\ii^{m+1})=(-1)^{(m+1)/2}$ gives
\eqref{eq:first-defect}.  Every other term in that equation is rational by
\cref{cor:finite-comb,thm:moment-egf}; solving for $D(m+1)$ proves
\eqref{eq:D-rational}.
\end{proof}

The first two evaluations are
\begin{equation}\label{eq:D2D4}
 D(2)=\frac{\pi^2}{18},
 \qquad
 D(4)=\frac{23\pi^4}{4050}.
\end{equation}
The first follows from the level-$1$ defect
$S_{1,1}-(I_1+1/4+1/48)=-1/144$.  For $m=3$ one obtains
\[
 S_{3,3}=\frac{42751}{147456},
 \qquad
 S_{3,3}-\left(I_3+\frac1{16}+\frac1{256}-\frac1{491520}\right)
 =\frac{23}{22118400},
\]
which gives the second identity.  Further exact values generated from
\eqref{eq:first-defect} appear in \cref{tab:dirichlet}.

The identity
\[
 H(\pi q)=\frac{2(-1)^{(q-1)/2}}{\pi q}H(\pi q/2)
 \qquad(q\ \text{odd})
\]
connects $D(s)$ to a twisted half-scale series.  This suggests a hierarchy of
functional relations combining odd residue classes, dyadic dilation, and
Thue--Morse signs.

\section{Other comb conventions and Rvachev/Fabius transforms}

Let
\[
 \Pi_{m,p}=h\sum_{k=0}^{N}(kh)^p
 =\frac{B_{p+1}(N+1)-B_{p+1}(0)}{(p+1)N^{p+1}}.
\]
The basic transformations require no new sampling:
\begin{align}
 h\sum_{k=0}^{N}(kh)^p\bigl(1-F(kh)\bigr)
 &=\Pi_{m,p}-S_{m,p},\label{eq:one-minus-F}\\
 h\sum_{k=0}^{N}(kh)^pF(1-kh)
 &=\Pi_{m,p}-S_{m,p}.\label{eq:F-reflected}
\end{align}
Because $u(x)=1-F(x)$ on $[0,1]$, these are also the positive-half Rvachev
sums.

The closed right rule, trapezoidal rule, and open/left rule satisfy
\begin{equation}\label{eq:rule-conversions}
 T_{m,p}=S_{m,p}-\frac h2,
 \qquad
 L_{m,p}=S_{m,p}-h.
\end{equation}
For the midpoint comb
\[
 M_{m,p}=h\sum_{k=0}^{N-1}
 \left(\frac{k+1/2}{N}\right)^p
 F\!\left(\frac{k+1/2}{N}\right),
\]
splitting the level-$(m+1)$ grid into even and odd nodes gives the exact
identity
\begin{equation}\label{eq:midpoint}
 M_{m,p}=2S_{m+1,p}-S_{m,p}.
\end{equation}
Thus every standard dyadic rule inherits the closed forms above.

For the full bump, \cref{thm:up-exactness} is the strongest statement:
\[
 h\sum_{k=-N}^{N}P(kh)u(kh)=\mathbb E P(X)
 \qquad(\deg P\le m).
\]
Translations and affine rescalings of $u$ merely transform the polynomial and
the mesh.  Tensor products therefore give an exact cubature rule on all
coordinatewise degrees at most $m$ for products of Rvachev bumps.

\section{Negative, fractional, and complex powers}

\subsection{Entire Mellin dependence}

For $\alpha\in\C$ define
\begin{equation}\label{eq:S-alpha}
 S_m(\alpha)=h\sum_{k=1}^{N}(kh)^\alpha F(kh),
 \qquad
 I(\alpha)=\int_0^1x^\alpha F(x)\dd x,
\end{equation}
using the real logarithm on $(0,1]$.  At $x=0$ the product is defined by its
limit $0$.

\begin{theorem}[Entire power parameter]\label{thm:entire-alpha}
Both $S_m(\alpha)$ and $I(\alpha)$ are entire functions of $\alpha$.  Moreover,
\begin{equation}\label{eq:fractional-finite}
 S_m(\alpha)=
 \frac{2^{-m(m+1)/2}}{m!N^{\alpha+1}}
 \sum_{k=1}^{N}k^\alpha
 \sum_{a=0}^{k-1}\tm(a)P_m\bigl(2(k-a)-1\bigr)
\end{equation}
is an exact finite expression for arbitrary complex, fractional, or negative
$\alpha$.
\end{theorem}

\begin{proof}
The finite sum is manifestly entire.  For $I(\alpha)$, fix a compact set
$K\subset\C$.  By flatness \eqref{eq:flat}, choose $A$ larger than
$-\min_{\alpha\in K}\Re\alpha$ plus any prescribed number of logarithmic
derivatives.  Then $x^\alpha(\log x)^rF(x)$ is uniformly integrable for
$\alpha\in K$.  Differentiation under the integral is valid to every order.
The finite formula follows from \cref{thm:dyadic-values}.
\end{proof}

The continuous entire function has a rapidly convergent Newton--binomial
series.

\begin{proposition}[Fractional moment series]\label{prop:fractional-series}
For every $\alpha\in\C$,
\begin{equation}\label{eq:I-alpha-series}
 I(\alpha)=\sum_{r=0}^{\infty}(-1)^r\binom{\alpha}{r}a_r,
\end{equation}
where $a_r$ is given explicitly by \eqref{eq:a-r}.  The convergence is locally
uniform in $\alpha$.
\end{proposition}

\begin{proof}
Use $F(x)=u(1-x)$ and expand $(1-y)^\alpha$ on $0\le y<1$.  Since $u$ is flat
at $y=1$, integration by parts or a beta estimate gives
$a_r=O_A(r^{-A})$ for every $A>0$.  On compact $\alpha$-sets the generalized
binomial coefficient grows at most polynomially in $r$, so the series is
normally convergent.
\end{proof}

For $\alpha=p\in\N_0$, the series terminates and reduces to
\eqref{eq:I-closed}.  For a negative integer it remains infinite but rapidly
convergent.  There is no singular transition at $\alpha=-1,-2,\ldots$; the
flat Fabius endpoint removes the Mellin poles that an ordinary nonzero density
would have.

\subsection{Universal endpoint expansion}

\begin{theorem}[Complex-power Euler--Maclaurin expansion]\label{thm:alpha-EM}
For every compact $K\subset\C$ and every fixed $R\ge1$, uniformly for
$\alpha\in K$,
\begin{equation}\label{eq:alpha-EM}
 S_m(\alpha)=I(\alpha)+\frac h2
 +\sum_{q=1}^{R-1}\frac{B_{2q}}{(2q)!}
   \fall{\alpha}{2q-1}h^{2q}
 +O_{K,R}(h^{2R})
 \qquad(m\to\infty).
\end{equation}
\end{theorem}

\begin{proof}
The function $x^\alpha F(x)$ extends smoothly and flatly through $x=0$ for
all $\alpha$ in a compact set.  At $x=1$, every positive derivative of $F$
vanishes and $F(1)=1$, hence
\[
 \left.\frac{\dd^r}{\dd x^r}\bigl(x^\alpha F(x)\bigr)\right|_{x=1}
 =\fall\alpha r.
\]
Apply Euler--Maclaurin to the right rule.  All lower-end corrections vanish.
\end{proof}

For natural $p$, the falling factorial becomes zero beyond degree $p$; at a
dyadic level above the alias threshold, \cref{thm:superconvergence} upgrades
this truncated asymptotic formula to an exact equality.  For nonintegral
$\alpha$ the series does not truncate.  Subtracting $h/2$ and then the $h^2$
term exposes the expected even-power convergence, as shown in
\cref{tab:fractional,fig:fractional-convergence}.

For a fixed negative power $\alpha=-s$, the first grid point is not dangerous:
\begin{equation}\label{eq:first-node}
 h^{1-s}F(h)=
 2^{m(s-1)}\frac{2^{-m(m+1)/2}}{m!}P_m(1),
\end{equation}
which decreases faster than every geometric sequence in $m$ for fixed $s$.
This gives a concrete discrete counterpart of Mellin entire-ness.

\section{Iterated sums of arbitrary integer order}

Let $a_k=(kh)^pF(kh)$ and define the inclusive discrete integral
\[
 (\J_h a)_K=h\sum_{j=0}^{K}a_j,
 \qquad
 \J_h^n=\underbrace{\J_h\circ\cdots\circ\J_h}_{n\ \mathrm{times}}.
\]

\begin{theorem}[Hockey-stick kernel]\label{thm:iterated-kernel}
For $n\ge1$,
\begin{equation}\label{eq:iterated-kernel}
 (\J_h^n a)_K
 =h^n\sum_{k=0}^{K}\binom{K-k+n-1}{n-1}a_k.
\end{equation}
At $K=N$ this becomes
\begin{equation}\label{eq:iterated-product}
 (\J_h^n[x^pF(x)])_N
 =\frac1{(n-1)!}\sum_{j=0}^{n-1}c_{n,j}(h)S_{m,p+j},
\end{equation}
where
\begin{equation}\label{eq:cnj}
 \prod_{r=1}^{n-1}(1+rh-x)
 =\sum_{j=0}^{n-1}c_{n,j}(h)x^j,
 \qquad
 c_{n,j}(h)=(-1)^j
 e_{n-1-j}(1+h,1+2h,\ldots,1+(n-1)h).
\end{equation}
\end{theorem}

\begin{proof}
Repeatedly interchange the finite sums; the number of weak chains between a
fixed $k$ and $K$ is the binomial coefficient in \eqref{eq:iterated-kernel}.
At $K=N$,
\[
 h^n\binom{N-k+n-1}{n-1}
 =\frac h{(n-1)!}\prod_{r=1}^{n-1}(1+rh-kh).
\]
Expand the polynomial and identify the ordinary comb moments.
\end{proof}

This answers the $n$th-order question without introducing a new family of
unknown constants.  Combining \cref{thm:iterated-kernel,thm:superconvergence}
gives the explicit stabilized rule
\begin{align}
 (\J_h^n[x^pF(x)])_N
 &=\frac1{(n-1)!}\sum_{j=0}^{n-1}c_{n,j}(h)
 \Biggl[
 I_{p+j}+\frac h2\notag\\
 &\hspace{28mm}+
 \sum_{q=1}^{\lfloor(p+j+1)/2\rfloor}
 \frac{B_{2q}}{(2q)!}\fall{p+j}{2q-1}h^{2q}
 \Biggr],\label{eq:iterated-stabilized}
\end{align}
valid, conservatively, for $m\ge p+n$.  It is an exact polynomial in $h$ with
rational coefficients.

The continuum limit is the usual $n$-fold integral:
\begin{equation}\label{eq:iterated-limit}
 \lim_{m\to\infty}(\J_h^n[x^pF(x)])_N
 =\frac1{(n-1)!}\int_0^1(1-x)^{n-1}x^pF(x)\dd x
 =\frac1{(n-1)!}\sum_{j=0}^{n-1}(-1)^j
   \binom{n-1}{j}I_{p+j}.
\end{equation}
For integral $p,n$, this limit is rational by \cref{thm:moment-egf}.

A strict nesting convention, in which every outer index is larger rather than
larger-or-equal, replaces the kernel by $\binom{K-k}{n-1}$ and changes
$1+rh-x$ to $1-rh-x$ (after the corresponding endpoint adjustment).  The same
polynomial reduction applies.

\subsection{Fractional-order repeated sums}

For $\Re\nu>0$ the standard discrete Riemann--Liouville continuation is
\begin{equation}\label{eq:fractional-iterate}
 (\J_h^\nu a)_N
 =h^\nu\sum_{k=0}^{N}
 \frac{\Gamma(N-k+\nu)}{\Gamma(\nu)\Gamma(N-k+1)}a_k.
\end{equation}
At $\nu=n\in\N$ it agrees with \eqref{eq:iterated-kernel}.  Combining it with
\eqref{eq:fractional-finite} gives a finite exact expression for arbitrary
complex power $\alpha$ and fractional order $\nu$.  Under ordinary compactness
conditions its continuum limit is
\begin{equation}\label{eq:fractional-limit}
 \frac1{\Gamma(\nu)}\int_0^1(1-x)^{\nu-1}x^\alpha F(x)\dd x.
\end{equation}
Unlike the integer-order kernel, the fractional kernel is not a polynomial, so
one should not expect a finite reduction to ordinary moments; a Newton-series
reduction is still available.

\section{Exact and numerical checks}

All exact tables in this section are generated by the accompanying Python
program from rational Bernoulli numbers, cumulants, Bell recurrences, and the
sparse Thue--Morse differences.  No floating-point values enter the proofs.
The fractional-power table uses 90--100 decimal digits and a level-$15$
endpoint-corrected reference.

% ed.: generated_results.tex was not shipped in the archive (the
% README lists the generated numerical material as part of the
% intended package).  This stub stands in so the document compiles.
\begin{center}
\fbox{\parbox{0.9\linewidth}{\itshape ed.: the generated numerical
results file (\texttt{generated\_results.tex}) was not shipped in the
archive; the tables and closed forms it would insert here are
unrecoverable.  The companion package
\texttt{Fabius\_Dyadic\_Comb\_Sums\_Report\_Package/} in this group
carries its own \texttt{generated/} data tables for the same circle
of identities.}}
\end{center}


\section{Further deductions}

\subsection{A quadrature-design interpretation}

The exactness in \cref{thm:up-exactness} is a Strang--Fix-type phenomenon in an
unusual guise.  Ordinarily, polynomial reproduction is encoded by zeros of a
Fourier transform at a reciprocal lattice.  Here the reciprocal-lattice zeros
are not merely present: their multiplicity increases linearly with the dyadic
level because each refinement introduces another sinc factor that vanishes at
all aliases.  The degree of exactness therefore grows with refinement while
the node weights remain the samples $h\,u(kh)$ themselves.

For the Fabius half-bump, cutting at the center removes one zero through the
factor $(\e^{\ii z}-1)^{-1}$.  The only residual aliases below degree $m$ are
the elementary endpoint pole, whose zeta sums are exactly the Bernoulli
corrections.  This gives a conceptual reason why the finite formula
\eqref{eq:stabilized} looks like Euler--Maclaurin but is stronger than it.

\subsection{Two-adic filtration of higher defects}

For $p\ge m$, write $\ell=2^rq$ with $q$ odd.  The local zero order of $R$ is
$m+r$, so only $0\le r\le p-m$ occur.  Hence
\[
 \Rema_{m,p}=\sum_{r=0}^{p-m}\Rema_{m,p}^{[r]},
\]
where each stratum is an odd-integer Dirichlet series built from a finite jet
of the zero-stripped sinc product at $\pi q$.  The $r=0$, $p=m$ jet is exactly
$D(m+1)$.  Higher jets involve $H^{(j)}(\pi q)$ and finite cotangent sums.  This
filtration is finite for every fixed pair $(m,p)$ and should be preferable to
directly differentiating the whole infinite product.

\subsection{Richardson extrapolation}

For complex $\alpha$, define $\widetilde S_m(\alpha)=S_m(\alpha)-2^{-m-1}$.
Then \eqref{eq:alpha-EM} contains only even powers $2^{-2m}$.  Standard
Richardson or Romberg elimination therefore applies with ratio $4$.  When
$\alpha=p\in\N_0$, the formal correction sequence is finite; once
$m\ge p+1$, eliminating all nonzero powers returns $I_p$ exactly in rational
arithmetic.  For nonintegral $\alpha$, the same scheme gives rapidly improving
approximants to the entire Mellin moment $I(\alpha)$.

\section{Conjectures}

The following claims are intentionally separated from the proved results.
They are supported by the exact experiments bundled with the report and by
the local alias structure, but a complete proof was not obtained here.

\begin{conjecture}[Sharp first nonzero defect]\label{conj:sharpness}
For every $m\ge1$, the smallest $p$ for which $\Rema_{m,p}\ne0$ is
\[
 p=\begin{cases}
 m,&m\ \text{odd},\\
 m+1,&m\ \text{even}.
 \end{cases}
\]
Equivalently, apart from the forced even boundary cancellation in
\cref{thm:superconvergence}, there are no accidental cancellations among the
odd-alias jets.
\end{conjecture}

\begin{conjecture}[Spectral positivity]\label{conj:positivity}
The rational numbers
\[
 d_r:=\frac{D(2r)}{\pi^{2r}}
\]
are positive for every $r\ge1$.
\end{conjecture}

The summands in $D(2r)$ are not all positive, so this is not termwise
trivial.  The $q=1$ term dominates for large $r$; the difficult part is a
uniform estimate covering the small exponents.

\begin{conjecture}[Arithmetic closure of higher alias jets]\label{conj:jets}
After grouping by $2$-adic valuation and by Fourier parity, every
parity-compatible higher-jet series needed for an integer comb defect reduces
to a rational multiple of the appropriate power of $\pi$.  More concretely,
the combinations of
\[
 \sum_{q\ \mathrm{odd}}\frac{H^{(j)}(\pi q)}{q^s}
\]
that occur in $\Rema_{m,p}$ belong to $\mathbb Q\pi^{s-j}$ whenever the total
parity makes the comb sum real.
\end{conjecture}

Rationality of the \emph{total} defect follows from the finite formula; the
content of the conjecture is that the valuation/jet decomposition can be made
arithmetically diagonal rather than relying on cancellation between distinct
series.

\section{Promising research directions}

\paragraph{Closed recurrences for $D(s)$.}
Use $H(z)=\sinc(z/2)H(z/2)$ to split odd residues modulo $2^r$ and seek a
Mahler-type functional equation for the entire function $D(s)$.  Exact values
at positive even integers provide unusually rigid interpolation data.

\paragraph{Explicit higher-defect algebra.}
Develop the zero-stripped Taylor coefficients in \eqref{eq:local-R} through
Bell polynomials of logarithmic derivatives.  Grouping by $v_2(\ell)$ should
produce a finite symbolic algorithm for $\Rema_{m,p}$ that is substantially
smaller than the direct dyadic-value sum when $p-m$ is small.

\paragraph{$q$-binomial compression.}
The polynomial $\Theta_m(z)$ is the finite Thue--Morse product, while
$P_m$ is controlled by the dyadic product $M(t)$.  A mixed generating function
in $m$, $z$, and the moment index may reveal a direct bridge to the
$q$-binomial identities already developed in the repository, potentially
compressing \eqref{eq:B-polylog} further.

\paragraph{Uniform complex-$\alpha$ error bounds.}
Combine the endpoint expansion \eqref{eq:alpha-EM} with the known
Lambert-$W$ asymptotics of $F$ near zero to obtain optimally truncated bounds
uniform in vertical strips of $\alpha$.  The key technical question is how the
flat but nonanalytic endpoint controls the beyond-all-orders remainder.

\paragraph{Fractional discrete calculus.}
Study the two-parameter function in \eqref{eq:fractional-iterate} jointly in
$(\alpha,\nu)$.  Its poles and zeros under analytic continuation in $\nu$ may
encode a discrete analogue of beta-moment identities for the inverse Fabius
function.

\paragraph{Multivariate cubature.}
Tensor products of \cref{thm:up-exactness} immediately give compactly
supported exact cubatures.  Non-tensor dyadic refinements, anisotropic sinc
products, and sparse grids are natural next steps.

\paragraph{Formal verification.}
The rational moment recurrence, sparse-difference generation, Prouhet
cancellation, and finite iterated-sum reduction are well suited to Lean.  The
analytic core can be split into reusable lemmas about infinite products,
Poisson summation, zero multiplicity, and removable singularities.

\section{Conclusion}

The weighted dyadic comb is much more rigid than a generic Riemann sum.  At an
arbitrary level, \cref{cor:finite-comb,thm:Gm} give a finite rational answer in
terms of Thue--Morse signs and Appell moment polynomials.  Above a simple
degree threshold, Poisson summation collapses that answer to a continuous
moment plus finitely many universal Bernoulli corrections.  The mechanism is
the growing multiplicity of the sinc-product zeros at dyadic aliases.

That same mechanism yields exact spectral Dirichlet values and a finite
$2$-adic filtration of every higher defect.  Flatness at the Fabius endpoint
makes negative and fractional powers analytically benign, while the integer
case acquires exact finite termination.  Finally, repeated summation introduces
no essentially new integer-order constants: its binomial kernel reduces it to
ordinary comb moments, and its limit is a rational beta-weighted moment.

\appendix
\section{Recursive repository inventory}\label{app:corpus}

The table below was generated from every \texttt{.tex} file below the requested
repository path, including all subdirectories.  Tags are lexical aids for the
audit, not semantic classifications.  The companion
\texttt{corpus\_audit.txt} records file sizes, first parsed headings, keyword
hit counts, the commit hash, and a combined SHA-256 digest.

% ed.: corpus_inventory.tex was not shipped in the archive.  This
% stub stands in so the document compiles.
\begin{center}
\fbox{\parbox{0.9\linewidth}{\itshape ed.: the relative-path corpus
inventory (\texttt{corpus\_inventory.tex}) was not shipped in the
archive; the audited file listing this appendix would reproduce is
unrecoverable.}}
\end{center}


\section{Reproducibility notes}

Run
\begin{lstlisting}[language=bash,basicstyle=\ttfamily\small,breaklines=true]
python experiments.py \
  --repo-docs /path/to/ProveIt/Analysis/FabiusFunction/docs \
  --out .
pdflatex fabius-dyadic-comb-sums.tex
pdflatex fabius-dyadic-comb-sums.tex
\end{lstlisting}
The program uses exact \texttt{fractions.Fraction} arithmetic for every
integer-power theorem and table.  High-precision fractional-power experiments
use \texttt{mpmath}; the optional figure uses \texttt{matplotlib}.  If those
optional packages are absent, the exact outputs are still generated.

\begin{thebibliography}{9}

\bibitem{Fabius1966}
J.~Fabius,
\newblock A probabilistic example of a nowhere analytic $C^\infty$-function,
\newblock \emph{Zeitschrift f\"ur Wahrscheinlichkeitstheorie und Verwandte
Gebiete} \textbf{5} (1966), 173--174.

\bibitem{Rvachev1990}
V.~L. Rvachev,
\newblock Compactly supported solutions of functional-differential equations
and their applications,
\newblock \emph{Russian Mathematical Surveys} \textbf{45}:1 (1990), 87--120.

\bibitem{GKP}
R.~L. Graham, D.~E. Knuth, and O.~Patashnik,
\newblock \emph{Concrete Mathematics}, 2nd ed.,
\newblock Addison--Wesley, 1994.

\bibitem{SteinShakarchi}
E.~M. Stein and R.~Shakarchi,
\newblock \emph{Fourier Analysis: An Introduction},
\newblock Princeton University Press, 2003.

\bibitem{ProveIt}
V.~Reshetnikov,
\newblock \emph{ProveIt: Analysis/FabiusFunction documentation},
\newblock recursive source corpus at commit \CorpusCommit, audited in this
report.

\end{thebibliography}

\end{document}
